\section*{Chapitre \ref{ch:g12:rlc-oscillations} --- Oscillations électriques libres~: RLC}

\begin{solution}{exo:g12:rlc-oscillations:1}
(a) $\sqrt{LC} = \qty{1.0e-4}{s}$~: $T_0 = \qty{0.63}{ms}$,
$f_0 \approx \qty{1.6}{kHz}$. (b) $\sqrt{LC} = \qty{3.16e-7}{s}$~:
$T_0 = \qty{2.0}{\micro s}$, $f_0 \approx \qty{0.50}{MHz}$ --- un
facteur $\sqrt{10^5} \approx 316$.
\end{solution}

\begin{solution}{exo:g12:rlc-oscillations:2}
$LC$~: $(\unit{V.s/A})(\unit{C/V}) = \unit{s} \times \unit{C/A} =
\unit{s} \times \unit{s} = \unit{s^2}$, donc $\sqrt{LC}$ est en
secondes. $2\pi$ est un nombre pur --- aucune dimension à changer.
\end{solution}

\begin{solution}{exo:g12:rlc-oscillations:3}
$Q_0 = CU_0 = \qty{2.8}{\micro C}$~; $E = \tfrac12 CU_0^2 =
\tfrac12 \times \num{4.7e-7} \times 36 \approx \qty{8.5}{\micro J}$.
La charge oscille à $f_0 = 1/(2\pi\sqrt{LC})$, le courant d'un quart
de \omterm{def:g12:oscillators-and-time:oscillator}{période} en décalage de
\omterm{def:g12:oscillators-and-time:phase}{phase}, l'\omterm{def:g10:energy-conservation:energy}{énergie} basculant entre
plaques et \omterm{def:g11:magnetic-fields:solenoid}{bobine}, total constant.
\end{solution}

\begin{solution}{exo:g12:rlc-oscillations:4}
$Q_0 = \qty{2.0}{\micro C}$~; $f_0 = 1/T_0 = \qty{250}{Hz}$~;
$q(T_0/2) = \qty{-2.0}{\micro C}$~; $i(0) = 0$ (cosinus plat en
$t=0$)~; $I_0 = 2\pi \times \num{2.0e-6}/\num{4.0e-3} \approx
\qty{3.1}{mA}$.
\end{solution}

\begin{solution}{exo:g12:rlc-oscillations:5}
Tout dans le \omterm{def:g12:rc-rl-circuits:capacitor}{condensateur} lorsque $q = \pm Q_0$,
$i = 0$ ($t = 0$, $T_0/2$, $T_0$, \dots)~: les extrêmes du pendule.
Tout dans la \omterm{def:g11:magnetic-fields:solenoid}{bobine} lorsque $q = 0$, $i = \pm I_0$
($t = T_0/4$, $3T_0/4$, \dots)~: le point le plus bas et le plus
rapide du pendule.
\end{solution}

\begin{solution}{exo:g12:rlc-oscillations:6}
$d^2q/dt^2 = -(2\pi/T_0)^2\,Q_0\cos(2\pi t/T_0) = -(2\pi/T_0)^2 q$,
donc $L\,d^2q/dt^2 + q/C = q\,[\,1/C - L(2\pi/T_0)^2\,] = 0$ pour
tout $t$ exactement lorsque $(2\pi/T_0)^2 = 1/(LC)$, c'est-à-dire
$T_0 = 2\pi\sqrt{LC}$.
\end{solution}

\begin{solution}{exo:g12:rlc-oscillations:7}
$L = \dfrac{1}{4\pi^2 f_0^2 C} = \dfrac{1}{39.5 \times
(\num{1.98e5})^2 \times \num{2.2e-10}} \approx \qty{2.9}{mH}$.
\end{solution}

\begin{solution}{exo:g12:rlc-oscillations:8}
$E = \tfrac12 CU_0^2 = \qty{5.0e-5}{J}$~; $\tfrac12 LI_0^2 = E$ donne
$I_0 = \sqrt{2E/L} = \sqrt{\num{1.0e-2}} = \qty{0.10}{A}$, atteint
pour la première fois à $T_0/4 = \tfrac14 \times \qty{0.63}{ms}
\approx \qty{0.16}{ms}$.
\end{solution}

\begin{solution}{exo:g12:rlc-oscillations:9}
$E_C \propto \cos^2(2\pi t/T_0)$, et $\cos^2$ se répète à chaque
\emph{demi}-tour~: deux remplissages du
\omterm{def:g12:rc-rl-circuits:capacitor}{condensateur} par \omterm{def:g12:oscillators-and-time:oscillator}{période}, d'où
$2f_0$. À $t = T_0/8$ la \omterm{def:g12:oscillators-and-time:phase}{phase} est $\pi/4$~:
$\cos^2(\pi/4) = \tfrac12$, donc $E_C = E_L = E/2$.
\end{solution}

\begin{solution}{exo:g12:rlc-oscillations:10}
Rapport $3.0/4.0 = 0.75$ par
\omterm{def:g12:rlc-oscillations:regimes}{pseudo-période}~: sixième maximum
$4.0 \times 0.75^6 \approx \qty{0.71}{V}$.
\omterm{def:g10:energy-conservation:energy}{Énergie} $\propto$
\omterm{def:g12:oscillators-and-time:oscillator}{amplitude} au carré~: $0.75^2 \approx 56\%$
survit à chaque \omterm{def:g12:oscillators-and-time:oscillator}{période}. La
\omterm{def:g12:rlc-oscillations:regimes}{pseudo-période} reste $\approx T_0 = 2\pi\sqrt{LC}$~:
un \omterm{def:g12:oscillators-and-time:damping}{amortissement} léger rétrécit
l'\omterm{def:g12:oscillators-and-time:oscillator}{amplitude}, à peine le tempo.
\end{solution}

\begin{solution}{exo:g12:rlc-oscillations:11}
$T_0 = 2\pi\sqrt{0.50 \times \num{2.0e-5}} = 2\pi \times
\num{3.16e-3} \approx \qty{20}{ms}$. Dictionnaire~: $k = m/(LC) =
0.50/\num{1.0e-5} = \qty{5.0e4}{N/m}$~; vérification
$2\pi\sqrt{m/k} = 2\pi\sqrt{\num{1.0e-5}}$ --- les mêmes \qty{20}{ms}.
\end{solution}

\begin{solution}{exo:g12:rlc-oscillations:12}
$C = \qty{500}{pF}$~: $f_0 = 1/(2\pi\sqrt{\num{1.0e-13}}) \approx
\qty{0.50}{MHz}$~; $C = \qty{50}{pF}$~: $f_0 \approx \qty{1.6}{MHz}$.
À $L$ fixée, $f_0 = (2\pi\sqrt{L})^{-1} C^{-1/2} \propto 1/\sqrt{C}$~:
$C$ multiplié par dix, seulement $\sqrt{10} \approx 3.2$ en $f_0$ ---
comme trouvé.
\end{solution}

\begin{solution}{exo:g12:rlc-oscillations:13}
$0.95^n = \tfrac12$~: $n = \ln 2/\lvert\ln 0.95\rvert \approx 13.5$,
donc environ $14$ \omterm{def:g12:oscillators-and-time:oscillator}{périodes}.
\omterm{def:g12:oscillators-and-time:oscillator}{Amplitude} $\propto \sqrt{E}$~: descend à
$1/\sqrt{2} \approx 0.71$ du départ. À \qty{1.0}{MHz},
$T_0 = \qty{1.0}{\micro s}$~: environ \qty{14}{\micro s}.
\end{solution}

\begin{solution}{exo:g12:rlc-oscillations:14}
Perte par cycle $0.020 \times \qty{1.0e-9}{J} = \qty{2.0e-11}{J}$, à
\num{32768}~cycles par seconde~: $P = \num{32768} \times \num{2.0e-11}
\approx \qty{6.6e-7}{W}$ --- moins d'un microwatt, des années sur
une pile bouton. En phase, parce que seule une poussée à la bonne
\omterm{def:g12:oscillators-and-time:phase}{phase} \emph{ajoute} de
l'\omterm{def:g10:energy-conservation:energy}{énergie}
(\omterm{def:g12:oscillators-and-time:resonance}{résonance})~; hors phase elle freinerait le
balancement.
\end{solution}

\begin{solution}{exo:g12:rlc-oscillations:15}
$C = \dfrac{1}{4\pi^2 f_0^2 L} = \dfrac{1}{39.5 \times
(\num{2.5e7})^2 \times \num{1.0e-6}} \approx \qty{41}{pF}$. Le
résistor doit pousser la boucle dans le
\omterm{def:g12:rlc-oscillations:regimes}{régime apériodique}~: la charge se fixe alors
sans dépasser --- pas d'oscillation, pas de sonnement.
\end{solution}

\begin{solution}{pb:g12:rlc-oscillations:1}
\textbf{1.} Loi de maille~: $u_L + u_C = 0$, avec $u_C = q/C$,
$u_L = L\,di/dt$, $i = dq/dt$~: $L\,d^2q/dt^2 + q/C = 0$.

\textbf{2.} $d^2q/dt^2 = -(2\pi/T_0)^2 q$, donc l'équation se lit
$q\,[\,1/C - L(2\pi/T_0)^2\,] = 0$, vrai pour tout $t$ exactement
lorsque $T_0 = 2\pi\sqrt{LC}$.

\textbf{3.} $LC$~: $(\unit{V.s/A})(\unit{C/V}) = \unit{s} \times
\unit{C/A} = \unit{s^2}$ --- $\sqrt{LC}$ est un temps.

\textbf{4.} $LC = \qty{6.6e-14}{s^2}$~: $T_0 = 2\pi \times \num{2.57e-7}
\approx \qty{1.6}{\micro s}$, $f_0 \approx \qty{620}{kHz}$ --- dans
\qtyrange{530}{1700}{kHz}.

\textbf{5.} $C = \dfrac{1}{4\pi^2 f_0^2 L} = \dfrac{1}{39.5 \times
\num{1.0e12} \times \num{2.0e-4}} \approx \qty{127}{pF}$.

\textbf{6.} Même formule à \qty{530}{kHz}~: $C \approx \qty{451}{pF}$.

\textbf{7.} À \qty{1700}{kHz}~: $C \approx \qty{44}{pF}$.

\textbf{8.} $C = \qty{480}{pF}$~: $f_0 \approx \qty{514}{kHz}$~;
$C = \qty{20}{pF}$~: $f_0 \approx \qty{2.5}{MHz}$ --- toute la bande,
avec de la marge aux deux extrémités.

\textbf{9.} $f_0 = (2\pi\sqrt{L})^{-1}C^{-1/2} \propto 1/\sqrt{C}$~;
doubler $f_0$ exige de diviser $C$ par $4$.

\textbf{10.} Puisque $f_0 \propto 1/\sqrt{C}$, des pas égaux de $C$
déplacent $f_0$ lentement aux grands $C$ et vite aux petits $C$~: un
écart donné entre stations prend alors peu de course de cadran ---
elles se serrent à l'extrémité des
\omterm{def:g12:oscillators-and-time:oscillator}{hautes fréquences} (faible $C$).

\textbf{11.} $Q_0 = CU_0 = \num{1.27e-10} \times 0.020 \approx
\qty{2.5e-12}{C}$~; $E = \tfrac12 CU_0^2 \approx \qty{2.5e-14}{J}$.

\textbf{12.} $\tfrac12 LI_0^2 = E$~: $I_0 = U_0\sqrt{C/L} =
0.020\sqrt{\num{6.35e-7}} \approx \qty{16}{\micro A}$.

\textbf{13.} $0.96^n = \tfrac12$~: $n = \ln 2/\lvert\ln 0.96\rvert
\approx 17$~cycles.

\textbf{14.} À \qty{1.00}{MHz}, $T_0 = \qty{1.0}{\micro s}$~: le
sonnement est divisé par deux en environ \qty{17}{\micro s}. Un
sonnement solitaire meurt en microsecondes~; la musique exige que la
boucle soit alimentée en continu.

\textbf{15.} L'\omterm{def:g10:signals-and-waves:wave}{onde} pousse des millions de fois par
seconde~; seulement lorsque sa \omterm{def:g12:oscillators-and-time:oscillator}{fréquence}
épouse $f_0$ les poussées arrivent en phase et construisent le
balancement --- la \omterm{def:g12:oscillators-and-time:resonance}{résonance}~; les stations
désaccordées poussent aussi souvent contre qu'avec. À l'émetteur, un
amplificateur renvoie l'\omterm{def:g10:energy-conservation:energy}{énergie}
dissipée chaque cycle~: un \omterm{def:g12:oscillators-and-time:oscillator}{oscillateur}
entretenu.

\textbf{16.} $LC = \qty{2.54e-14}{s^2}$~: $k = m/(LC) =
0.20/\num{2.54e-14} \approx \qty{7.9e12}{N/m}$.

\textbf{17.} Environ $10^8$~fois le ressort de voiture~: aucun
système masse--ressort de paillasse n'atteint les mégahertz --- les
raideurs et masses du quotidien sont désespérément lentes.

\textbf{18.} $m = \dfrac{k}{(2\pi f_0)^2} = \dfrac{100}
{(\num{6.28e6})^2} \approx \qty{2.5e-12}{kg}$ --- quelques nanogrammes,
une poussière.

\textbf{19.} Un cristal de quartz~: un lambeau ultra-léger et
ultra-raide chantant à sa \omterm{def:g12:rlc-oscillations:period}{fréquence propre} --- le
garde-temps rencontré plus tôt cette année.

\textbf{20.} Avec $L = \qty{0.20}{mH}$ et \qtyrange{20}{480}{pF} la
boucle balaie environ \qtyrange{0.51}{2.5}{MHz}, couvrant la bande~;
$C \approx \qty{127}{pF}$ choisit \qty{1.00}{MHz}~; laissé seul le
sonnement se divise par deux en $\approx \qty{17}{\micro s}$, mais
l'\omterm{def:g10:signals-and-waves:wave}{onde} de la station, résonnant en phase, réalimente
la boucle --- les nouvelles jouent sur
l'\omterm{def:g10:energy-conservation:energy}{énergie} de l'\omterm{def:g10:signals-and-waves:wave}{onde}
elle-même.
\end{solution}
